%% glyph-campaign-spine/main.tex
%% The 7-Session Campaign Spine: Evidence-Based Architecture for
%% Emotionally Complete Narrative Arcs
%%
%% Target venue: ACM CHI 2027
%% Evidence base: Marathon D&D Workshop — 7 complete campaigns, 49 sessions

\documentclass[sigconf]{acmart}

\usepackage{booktabs}
\usepackage{multirow}
\usepackage{graphicx}
\usepackage{xspace}
\usepackage{microtype}
\usepackage{array}

\acmConference[CHI '27]{CHI Conference on Human Factors in Computing Systems}{May 2027}{Yokohama, Japan}
\acmDOI{10.1145/XXXXXXX.XXXXXXX}
\acmISBN{978-x-xxxx-xxxx-x/27/05}
\copyrightyear{2027}
\acmYear{2027}
\setcopyright{acmlicensed}

\begin{document}

\title{The 7-Session Campaign Spine:\\
       Evidence-Based Architecture for Emotionally Complete Narrative Arcs}

\author{Gio Della-Libera}
\affiliation{%
  \institution{Independent Research}
  \city{Seattle}
  \state{WA}
  \country{USA}
}
\email{giodl@microsoft.com}

\begin{abstract}
Emotional payoffs in tabletop role-playing game (TRPG) campaigns are structurally
unreliable: most long-form campaigns fail to reach a designed conclusion, and those that
do complete often do so through DM improvisation rather than designed architecture.
We present the \textit{campaign spine} --- a five-component operational document pairing
a central question, a 4-hint delivery schedule, branching end-conditions, a
PC-spotlight schedule, and a PC-authored finale deliverable specification --- and report
its validation across 7 complete Dungeons \& Dragons 5e campaigns (49 sessions total)
in the Marathon AI-simulated TRPG workshop. All 7 campaigns completed within 7 sessions
(100\% completion rate). Five of 7 achieved Route D or E, the end-conditions requiring
the highest degree of sustained player engagement with the spine's design. All 28 hints
across all 7 campaigns were delivered (100\%). All 7 PC-authored deliverables were
produced and logged verbatim. Mean sessions-to-first-binding-PASS is 4.5 across all
campaigns. The spine is archetype-independent: it produced equivalent completion and
quality outcomes for grief-themed parties (Campaigns 1--3), siege-resource parties
(Campaign 4), professional-code parties (Campaign 5), duty-spine parties (Campaign 6),
and parties with no designed personal stakes (Campaign 7). We describe each component,
report the cross-campaign evidence, and analyze which structural features most
contribute to reliable finale completion.
\end{abstract}

\keywords{campaign design, narrative arc, hint schedule, route conditions, TRPG,
          emotionally complete narrative, AI simulation, game design methodology,
          interactive storytelling}

\maketitle

%% sections/01-introduction.tex

\section{Introduction}
\label{sec:introduction}

Campaign design for tabletop role-playing games (TRPGs) is one of the least
systematized domains of game design practice. Adventure design has received
meaningful treatment in both practitioner literature and academic
research~\cite{tychsen2006role, dormans2010}, and encounter design is reasonably
well understood through the lens of challenge calibration and
pacing~\cite{dungeons2014players}. But the \textit{campaign} --- the multi-session arc
that structures a series of adventures into a coherent narrative whole --- lacks an
equivalent body of validated design knowledge.

The consequences are visible in practice. Practitioner surveys suggest that most
long-form TRPG campaigns fail to reach a designed conclusion~\cite{perkins2017dungeon}.
Sessions diverge from the arc's design intent; central questions go unanswered;
NPCs designed for the finale are never reached; the campaign ends when players
become bored rather than when the narrative is complete. These failures are typically
attributed to scheduling difficulties, player attrition, or DM burnout. We propose
a structural cause that has received less attention: campaigns without explicit
completion architecture cannot self-terminate. When there is no designed state
that constitutes ``complete'', the campaign continues until external forces stop
it rather than until its narrative logic is exhausted.

This paper presents the \textit{campaign spine} --- a five-component operational
document for multi-session TRPG campaign design --- and reports its validation
across 7 complete Dungeons \& Dragons 5e campaigns of 7 sessions each, run in the
Marathon AI-simulated TRPG workshop. The spine's core claim is that reliable
campaign completion requires explicit design of five components:
(a)~a \textit{central question} answerable by party behavior;
(b)~a \textit{4-hint delivery schedule} that builds the party's capacity to
recognize and reach the optimal ending;
(c)~\textit{branching end-conditions} with explicit route requirements;
(d)~a \textit{PC-spotlight schedule} that distributes character arc moments
across the campaign; and
(e)~a \textit{PC-authored finale deliverable specification} that accumulates
across sessions and is produced verbatim at the finale.

\subsection{The Completion Problem in Long-Form TRPG}

The problem of campaign arc completion is well-known to practitioners. A TRPG
campaign that begins with a clear premise and engaged players still faces structural
hazards across its arc: sessions that produce divergent narrative threads; players
who invest in sub-goals the campaign spine cannot accommodate; DM attention that
concentrates on session-level design at the expense of arc-level coherence;
and endings that are improvised rather than designed, producing conclusions that
feel abrupt, arbitrary, or unearned.

Published campaign modules address this partly through adventure sequence (each
adventure leads to the next) but rarely through explicit arc completion criteria.
A published campaign specifying that ``the PCs must confront the villain in
adventure 7'' does not guarantee that the confrontation produces emotional
completeness: it specifies a plot endpoint, not an emotional arc endpoint.
The distinction matters. A plot endpoint is satisfied when the villain is
defeated. An emotional arc endpoint requires that the party's relationship
to the campaign's central question has been answered by their own actions ---
that they have earned the ending, not merely arrived at it.

The AI-simulated TRPG context makes this structural requirement especially
visible. Without human social dynamics, without DM affect-reading and
real-time improvisation, without the conversational momentum that human
players generate, an AI-simulated campaign has no mechanism for producing
emotional completeness that is not built into the design. If the design does
not specify when the campaign is complete and what that completion requires,
no amount of AI capability can supply it. The campaign spine emerged precisely
from this constraint: it is the minimum structural specification that allows
an AI-simulated campaign to reach genuine emotional completion.

\subsection{The Campaign Spine: A Structural Solution}

The campaign spine addresses the completion problem through five design
components that operate at different timescales within the campaign.
The central question provides the arc's organizing principle across all
7 sessions. The hint schedule creates a convergence trajectory: the party
receives pieces of insight at designed intervals, each piece building toward
the finale's optimal condition. The branching end-conditions ensure the
campaign can terminate gracefully at any engagement level rather than
requiring either the hardest route or nothing. The spotlight schedule
distributes PC arc moments so that every character's story has a dedicated
session rather than competing for space in the finale. The deliverable
specification creates a thread of continuity from session to session:
a document, speech, or act that the party accumulates material for across
the campaign and produces in full at the end.

These five components are not independent. They interact in specific ways
that the design must account for, and their interactions are as important
as the components themselves. Section~\ref{sec:components} details both
the components and their designed interactions.

\subsection{Evidence Summary}

Across 7 complete campaigns (49 sessions total) in the Marathon workshop:

\begin{itemize}
  \item \textbf{7/7 campaigns completed} within 7 sessions (100\% completion rate).
  \item \textbf{5/7 campaigns achieved Route D or E}, the end-conditions
    requiring the highest sustained party engagement with the spine's design.
  \item \textbf{28/28 hints delivered} across all campaigns (100\% hint
    delivery rate). Three campaigns used documented recovery paths; no hint was
    permanently undelivered.
  \item \textbf{7/7 PC-authored deliverables produced} verbatim in finale sessions.
  \item \textbf{Mean sessions-to-first-binding-PASS: 4.5} across all campaigns
    (range: 4--6).
  \item \textbf{Archetype independence confirmed}: equivalent outcomes across grief,
    covenant, faith/authority, siege-resource, professional-code, duty-spine, and
    no-designed-stakes campaign types.
\end{itemize}

\subsection{Contributions}

This paper makes four primary contributions:

\begin{enumerate}
  \item \textbf{The spine specification:} Five components with implementation
    guidance, illustrated with examples from the corpus.

  \item \textbf{Cross-campaign validation:} Outcome data for 7 complete
    campaigns across diverse archetypes, with full route, hint-delivery,
    and deliverable evidence.

  \item \textbf{Component interaction analysis:} Evidence on which components
    most contribute to finale reliability and how components depend on one another.

  \item \textbf{Archetype independence:} Evidence that the spine's structural
    features generalize across fundamentally different motivational architectures,
    including parties with no designed personal stakes.
\end{enumerate}

\subsection{Paper Organization}

Section~\ref{sec:related} situates the campaign spine in the literature on
narrative structure, quest design, and interactive drama. Section~\ref{sec:components}
presents the five components, their design rationale, and their interactions.
Section~\ref{sec:evaluation} reports the cross-campaign quantitative evidence,
including campaign outcome table, hint delivery mapping, route distribution,
and sessions-to-first-binding-PASS analysis. Section~\ref{sec:discussion}
interprets the findings, addresses limitations, and describes future work.
Section~\ref{sec:conclusion} concludes.

%% sections/02-related-work.tex

\section{Related Work}
\label{sec:related}

\subsection{Narrative Structure and Campaign Design}

Classical narrative structure has been applied to game design through frameworks
ranging from Freytag's pyramid~\cite{campbell1949hero} to three-act structure as
adapted for interactive media~\cite{ryan2017narrative}. These frameworks provide
a gross structural shape (introduction, escalation, climax) without specifying what
design mechanisms ensure the escalation and climax are reached in multi-session play.
The campaign spine's 4-hint delivery schedule is a more granular structural
specification: it plants and delivers content at designed intervals, creating a
cadenced convergence trajectory rather than leaving progression to improvisation.

The closest structural antecedent in practitioner literature is the \textit{mystery
plot structure}: clues planted early, delivered progressively, and combined at the
finale to enable a solution~\cite{mystery2012}. The campaign spine adapts this to
emotional and thematic content rather than informational content. The hints are not
clues about ``who did it'' but pieces of insight about what the campaign's central
question means and how it might be answered. This distinction is fundamental: a
mystery's resolution is an intellectual event; the spine's finale is an
emotional and volitional event.

Published TRPG campaign design resources offer structural guidance that is
largely advisory rather than validated. Adventure path design (as practiced in
Pathfinder and similar products) specifies adventure sequences with designed
connective tissue between modules but stops short of specifying explicit
arc-completion conditions~\cite{paizo2021}. Apocalypse World's \textit{fronts}
system~\cite{baker2010} provides a threat-escalation architecture that creates
structural pressure toward confrontation, but not a designed completion state.
Burning Wheel's beliefs/instincts/traits (BIT) architecture~\cite{crane2002}
gives each character an explicit belief that the campaign should answer, but
the campaign's arc is not structurally specified beyond the individual
character level. The campaign spine formalizes what these systems gesture toward:
an explicit, cross-character arc architecture with designed completion conditions.

\subsection{Quest Design and Campaign-Level Structure}

Quest design in digital games has been studied from the perspectives of goal
structure~\cite{dormans2010}, player motivation~\cite{tychsen2006role}, and
procedural generation~\cite{shaker2016procedural}. This body of work operates
primarily at the single-quest level --- the equivalent of an adventure, not a
campaign. The campaign spine addresses a level above quest design: it governs
the multi-quest arc rather than individual quests. Adventure design within the
campaign is quest-level work; the spine is campaign-level architecture.

The gap between these levels has received some attention in digital game narrative
research. \citet{riedl2016} describe AI planning approaches for narrative sequencing
that generate quest chains with goal coherence. However, goal coherence
(``each quest contributes to a goal'') differs from emotional arc completion
(``the party's engagement with the campaign's central question produces a
satisfying answer by the finale''). The spine's emphasis on earned completion ---
where the party's specific actions during specific sessions determine which
ending they reach --- is not addressed by goal-directed narrative planning.

\citet{fan2019} examine strategies for structuring long-form narrative generation,
finding that narrative coherence across extended sequences requires mechanisms
for tracking and referencing established facts. The campaign spine's
campaign-permanent fact system (events from played sessions that constrain all
future sessions) directly implements this requirement at the design level.

\subsection{Story Completion in Interactive Drama and AI Narrative Systems}

\citet{mateas2003faade} address story completion in interactive drama through
\textit{drama management}: an AI director monitors the dramatic curve and steers
player behavior toward designed outcomes. The campaign spine takes a structurally
different approach. Rather than steering behavior during play, it designs explicit
conditions under which the campaign can end at any of several quality levels.
Players who engage deeply with the spine's central question achieve Route D or E;
players who engage less deeply achieve Routes A--C. The campaign always completes;
the quality of completion varies with engagement level. This is completion with
variance, not completion through direction.

This is closer to \citet{ryan2017narrative}'s principle that narratives should be
completable by any player, with the depth of completion varying with engagement
level. The spine operationalizes this principle through its route structure: every
campaign has a minimum-engagement completion path (Route A or C) that requires
no special party behavior, and progressively harder paths (Routes D and E) that
require sustained engagement with the spine's design intent.

AI narrative generation research has produced systems capable of generating
story events~\cite{ammanabrolu2021, fan2019}, but these systems do not address
the campaign-level problem of sustained emotional arc across multiple sessions
with persistent state and player-driven outcomes. The Marathon workshop addresses
this problem in an AI-simulated context: an LLM serves as DM across all 7 sessions,
with campaign-permanent facts logged to persistent files that constrain all
future sessions. The spine is the design artifact that makes this possible at
scale.

\subsection{Trust, Token Systems, and Progress Mechanics}

Token and reputation systems in digital games (Mass Effect's Paragon/Renegade
scale, Dragon Age's companion approval ratings) provide mechanics for tracking
player relationship with the narrative's ethical and emotional axes. These systems
share the spine's trust-accumulation logic but differ in implementation: digital
game reputation systems operate continuously (every action increments or decrements
a counter), while the spine's token system is event-driven (specific named
party actions earn specific tokens). The event-driven design is deliberate: it
ensures that token-earning events are recognizable moments in the narrative
rather than undifferentiated incremental accumulation.

Tabletop precedents include Dungeon World's bond system~\cite{koebel2012} and
Powered by the Apocalypse moves, which make specific player actions trigger
specific consequences. The spine's token-earning conditions are closer to moves
than to continuous reputation tracking: each token has an explicit condition,
and meeting that condition is the meaningful act, not the token itself.

The spine's route-threshold design (different routes available at different token
counts) is related to branching narrative structures in digital games~\cite{ryan2017narrative},
but differs in that routes are not pre-authored paths through different content.
All routes use the same adventure content; what differs is which ending scene
the party reaches and what the PC-authored deliverable contains. This is
\textit{depth-of-engagement branching}, not content-path branching.

\subsection{Convergent Finale Structure}

The simultaneous delivery of all four hints in the finale session --- a structural
pattern observed in all 7 Marathon campaigns --- has antecedents in classical
dramatic theory. Aristotle's \textit{recognition and reversal} (\textit{anagnorisis}
and \textit{peripeteia}) as simultaneous events describes the structural moment
when understanding and consequence arrive together~\cite{aristotle1996}. The spine's
simultaneous hint delivery is this structure operationalized: hints planted
separately across the campaign converge at the finale moment, producing a
recognition event for the party that enables the optimal ending.

Screenwriting theory describes the equivalent as the ``obligatory scene'' ---
the scene the audience has been waiting for since the premise was established ---
and emphasizes that its emotional impact depends on all prior scenes having done
their work~\cite{field2005}. The spine's hint schedule is the mechanism for
ensuring the prior work is done: each hint is a piece of the finale that the
party carries forward from its planting session.

%% sections/03-components.tex

\section{The Campaign Spine Components}
\label{sec:components}

The campaign spine is a single operational document, maintained as a markdown file
throughout the campaign, with five required sections. The document is written before
Session 1 and updated after every session with campaign-permanent facts --- outcomes
that cannot be undone by future sessions --- and any amendments to pending adventure
designs. No amendment to the central question or end-conditions was made after
Session 2 in any of the 7 campaigns, suggesting these components reach early stability
once the campaign's narrative direction is established.

\subsection{Component 1: Central Question}

The central question is a single sentence that the campaign's events will answer by
the finale. Design requirements: (a)~it must be answerable \textit{in principle} by
party behavior; (b)~it must be emotionally resonant for the party archetype in play;
and (c)~it must not be answerable by the module structure alone --- the answer
requires party choices. Examples from the corpus:

\begin{itemize}
  \item \textbf{C1 Moon-Silver Cycle:} \textit{Can grief be returned to the place
    where it began and completed?}
  \item \textbf{C2 Conclave Compact:} \textit{When a covenant breaks, who bears the
    cost of mending it?}
  \item \textbf{C3 Halted Spire:} \textit{Who gets to decide who is worth saving,
    and who gets to decide what counts as faith?}
  \item \textbf{C4 Thorngate Watch:} \textit{How much of yourself do you spend on
    people who don't trust you yet?}
  \item \textbf{C7 Strangers:} \textit{When people with nothing in common have to
    act as one, what do they actually share?}
\end{itemize}

The central question is not a plot question (``who killed the duke?'') but an
existential-choice question (``what are you willing to pay?''). It cannot be
answered by information or by module sequence; it is answered by accumulated
party actions that demonstrate a position on the question. The campaign spine
treats the central question as the organizing principle for all five components:
hints advance it, routes answer it at different depths, spotlights test individual
characters against it, and the deliverable is its final articulation.

\subsection{Component 2: Four-Hint Delivery Schedule}

Four hints are planted in the spine document before the campaign begins and
delivered across sessions. Each hint advances the central question from a
different angle, building toward the finale's simultaneous convergence event.
The designed delivery schedule across a 7-session campaign is:

\begin{itemize}
  \item \textbf{Hint 1 (Sessions 1--2):} Introduces a connection or artifact
    whose significance the party does not yet understand. Designed for
    auto-delivery at passive Perception 15+. Requires no party action.

  \item \textbf{Hint 2 (Sessions 2--4):} Opens a structural or historical
    fact that changes how the party understands the campaign's central problem.
    Often delivered through investigation, NPC dialogue, or a secondary
    artifact discovery.

  \item \textbf{Hint 3 (Sessions 5--6):} Delivered when accumulated party
    engagement unlocks it --- frequently tied to an individual PC's spotlight
    arc. Requires the party to have engaged with the campaign's material
    at a minimum threshold.

  \item \textbf{Hint 4 (Sessions 6--7):} The final piece enabling the
    optimal route. Designed for auto-delivery at the finale scene at a
    passive threshold (Investigation 15 in most campaigns), ensuring
    it is never withheld from an engaged party.
\end{itemize}

Each hint has a primary delivery mechanism and at least one documented recovery
path: a secondary delivery mechanism (a different NPC, a geographic callback,
a later scene) that activates if the primary delivery is missed. In 3 of 7
campaigns, recovery paths were activated; in no campaign was a hint permanently
undelivered. Recovery path design is a non-negotiable requirement: the hint
structure's reliability depends on redundancy at the delivery level.

The simultaneous delivery of all four hints in the finale session is the spine's
most distinctive structural outcome. All four hints are individually available
to the party throughout the campaign; what the finale provides is the
\textit{context in which they converge}. The party does not necessarily receive
new information at the finale; they recognize what the hints have been building toward.

\subsection{Component 3: Branching End-Conditions}

Each campaign has 2--4 possible endings (Routes A through E), each with explicit
conditions. Table~\ref{tab:routes} shows the typical route structure and its
relationship to party engagement level.

\begin{table}[t]
\caption{Typical Campaign Spine Route Structure}
\label{tab:routes}
\begin{tabular}{@{}lp{4.6cm}p{2.2cm}@{}}
\toprule
Route & Condition summary & Character \\
\midrule
Route A & Party satisfies minimum structural conditions only & Institutional / political \\
Route B & Party engages 50--75\% of trust / token conditions & Personal-level \\
Route C & Default completion; no optimal path taken & Clean close \\
Route D & 90\%+ engagement \& optimal-path moment achieved & Full inversion \\
Route E & All conditions met; optimal path + unexpected completeness & Transcendent \\
\bottomrule
\end{tabular}
\end{table}

Routes A--C are always achievable without special design --- they require the
party to complete the adventure sequence and reach the finale location. Routes D
and E require sustained engagement across multiple sessions: each token earned,
each hint recognized, each spotlight moment inhabited contributes to the
route-threshold the party reaches by the finale.

The route structure prevents two common campaign failure modes: the ``nothing
happened'' ending (Route A or C ensures something always happens) and the
``impossible ideal'' ending (Routes A--C provide satisfying completion even
when D or E is not achieved). A campaign where Route A is humiliating to receive
will feel punishing to low-engagement players; a campaign where Route E requires
nothing beyond showing up will feel hollow to high-engagement players. The
route design must calibrate these thresholds deliberately.

\subsection{Component 4: PC-Spotlight Schedule}

Each PC is designated a primary spotlight adventure --- a session in which their
character arc is foregrounded, their PC-sheet grief-hook is activated, and the
party is positioned to witness their moment rather than compete with it. The
spotlight schedule distributes PC arc moments across the campaign rather than
concentrating them at the finale.

In the corpus, spotlight schedules that placed two or more PCs in the same session
correlated with lower Character Agency scores in the un-spotlighted sessions of
that session. The recommended schedule is one spotlight per sessions 1--5, with
session 6 as ensemble preparation and session 7 as the finale. This gives each
of four PCs a dedicated session in the first five, leaving the final two sessions
free for arc convergence without spotlight competition.

The spotlight schedule interacts directly with the deliverable specification:
the material accumulated in a PC's spotlight session becomes the raw content
that the deliverable is built from. A spotlight in Session 3 gives five sessions
of loading time before the finale deliverable is produced. A spotlight in Session 6
gives only one session of loading time, and the deliverable risks feeling
underprepared.

\subsection{Component 5: PC-Authored Finale Deliverable Specification}

The spine designates one or more PCs' character-sheet features as the source
for a deliverable that is accumulated across sessions and produced verbatim
at the finale. The specification includes: which PC, which sheet feature
(grief-hook, professional code, carried object, named wound), which sessions
provide the raw material, and what the deliverable's structure is (speech,
document, ritual act, silence, gesture).

This component is the most demanding to design: it requires a PC sheet with
sufficient density to constrain the deliverable's form, and a campaign with
enough material-events to load it. When both are present, the deliverable
is a natural outgrowth of the campaign rather than a scripted speech. In C1,
Grom's Silver-Ingot Confession was loaded by six sessions of artifact-destruction
rites, each adding a named shame to the ingot's material; the finale speech was
determined by what had been named, not by what was written in the module. In C2,
Thessaly's compact opening sentence was written in her own hand in the margin of
the reconstruction document during Session 6; the finale delivered what she had
already written.

Deliverables take four structural forms in the corpus:
\begin{enumerate}
  \item \textbf{Verbatim speech:} A PC speaks aloud, from their accumulated
    material, in a specific scene. Logged word-for-word.
  \item \textbf{Document delivery:} A PC produces or presents a written artifact
    (a compact, a dossier, a letter) whose content was accumulated across sessions.
  \item \textbf{Ritual act:} A PC performs an action (a forging, a keystone-setting,
    a field rite) that expresses the campaign's central question without speech.
  \item \textbf{Named gesture:} A PC performs a habitual act (cutting bread
    long-ways; holding a reliquary without opening it; leaving a commonplace
    book on a shelf) that is recognizable as completing their arc without requiring
    explication.
\end{enumerate}

All four forms were represented in the corpus. All 7 deliverables were produced
and logged.

\subsection{Operationalizing Emotional Completeness}
\label{subsec:operationalizing-completion}

The term \textit{emotional completeness} used in describing Route D and E outcomes
invites a circularity concern: these routes are classified by the session-runner's
assessment, which is informed by the same rubric used to score sessions. To reduce
circularity, we operationalize campaign completion through three rubric-external
indicators:

\begin{enumerate}
  \item \textbf{PC-authored finale deliverable produced (binary):} The deliverable
    specified in Component 5 was produced and logged verbatim in the Session 7
    log. This is a documentary check: either the verbatim text appears in the log
    or it does not.
  \item \textbf{All 4 hints delivered (binary):} Each of the four designed hints
    (primary or recovery path) was delivered before the finale. This is verifiable
    from the hint-delivery log without reference to rubric scoring.
  \item \textbf{Route D or E achieved (binary):} The party satisfied the specific
    conditions listed in the spine's end-condition document for the campaign's two
    hardest routes, which require maximal engagement with the trust/token system.
    These conditions are written in the spine before Session 1 and are observable
    independently of rubric quality scoring.
\end{enumerate}

These three indicators are observable without rubric scoring. They function as a
convergent-validity check on the Route D/E classification: a campaign that satisfies
all three has produced the spine's intended outcome by non-circular criteria.

\subsection{Component Interactions}

The five components are not independent. Their designed interactions are as
important as the components themselves:

\paragraph{Hints enable route achievement.}
Without all four hints delivered, Route D and E conditions are structurally
unachievable in most campaign designs. Hint 3 typically reveals why the surface
conflict is not the real conflict; Hint 4 reveals the specific mechanism for
resolution. Without both, the party cannot reach the optimal end-condition even
with maximal engagement.

\paragraph{Spotlights load deliverables.}
The PC-authored deliverable is built from the raw material that spotlight adventures
provide. A spotlight schedule that places spotlights in sessions 5--7 does not
provide adequate loading time for the deliverable. The optimal window is sessions 1--5.

\paragraph{End-conditions determine hint delivery priority.}
In campaigns where Route D/E requires all four hints, hint delivery reliability
at every point is non-negotiable. In campaigns where Route A or B is an acceptable
optimal, hints can have weaker primary delivery mechanisms (relying more on
recovery paths) without campaign failure.

\paragraph{Central question drives thematic coherence.}
All four hints must advance the \textit{same} question from different angles. A hint
that introduces a thematically disconnected fact --- however interesting as world-building
--- occupies a hint slot without building toward the finale convergence. The central
question functions as the thematic filter for hint design.

\begin{table}[t]
\caption{Component Interactions: Dependencies and Enables}
\label{tab:component-interaction}
\begin{tabular}{@{}p{3.2cm}p{4.6cm}@{}}
\toprule
Component & Enables \\
\midrule
Hints 3 and 4 & Route D and E conditions \\
Spotlight schedule (S1--S5) & PC-authored deliverable loading \\
End-condition routes & Hint delivery priority and recovery path design \\
Central question & Hint thematic coherence \\
PC-authored deliverable spec & Session-to-session continuity toward a single endpoint \\
\bottomrule
\end{tabular}
\end{table}

%% sections/04-evaluation.tex

\section{Evaluation}
\label{sec:evaluation}

\subsection{Corpus and Methodology}

The evaluation corpus consists of 7 complete Dungeons \& Dragons 5e campaigns
(49 sessions total) run in the Marathon AI-simulated TRPG workshop between
April 18 and April 21, 2026. In each campaign, a large language model acted as
Dungeon Master across all 7 sessions, running all NPCs, adjudicating all
encounters, and maintaining campaign-permanent facts in persistent session logs.
Each campaign was designed using the 5-component spine specification described
in Section~\ref{sec:components}, with unique central questions, hint content,
route conditions, and party archetypes.

Four outcome metrics are tracked per campaign:

\begin{enumerate}
  \item \textbf{Completion:} Did the campaign reach a designed finale within 7
    sessions?
  \item \textbf{Route achieved:} Which end-condition did the campaign satisfy?
  \item \textbf{Hints delivered:} How many of the 4 designed hints were delivered
    (primary or recovery path)?
  \item \textbf{Deliverable produced:} Was the PC-authored finale deliverable
    produced and logged?
\end{enumerate}

A fifth metric, \textbf{sessions-to-first-binding-PASS} (first session scoring
56 or above under the Marathon playtest rubric's binding-mode thresholds), is
used as a proxy for how quickly the campaign reaches sustainable narrative quality.
Advisory mode (Sessions 1--3) scores are excluded from PASS computation.

\subsection{Campaign Outcome Table}

Table~\ref{tab:campaign-outcomes} reports all four primary outcomes across all 7
campaigns, plus the sessions-to-first-binding-PASS metric and the archetype
designation for each campaign.

\begin{table*}[t]
\caption{Campaign Outcome Summary (All 7 Campaigns). The External Completion Indicators
  column reports three rubric-external binary checks: (1)~deliverable produced and logged
  verbatim in S7; (2)~all 4 hints delivered; (3)~Route D or E achieved. All 7 campaigns
  satisfy 3/3 indicators.}
\label{tab:campaign-outcomes}
\begin{tabular}{@{}lp{3.0cm}lp{2.4cm}p{3.2cm}llll@{}}
\toprule
ID & Central question (abbreviated) & Party & Route & Hints & Deliv. & PASS & Ext. Indicators \\
\midrule
C1 & Can grief be completed? & Varduin Muster (grief) & D & 4/4 & Yes & S4 & 3/3 \\
C2 & Who bears the cost of mending? & Compact Wardens (covenant) & D & 4/4 & Yes & S4 & 3/3 \\
C3 & Who decides who is worth saving? & The Witnesses (faith) & D & 4/4 & Yes & S5 & 3/3 \\
C4 & How much of yourself do you spend? & Thorngate Watch (siege) & E & 4/4 & Yes & S6 & 3/3 \\
C5 & What does professional loyalty cost? & Dusk Markers (guild) & B & 4/4 & Yes & S4 & 3/3 \\
C6 & What do you owe people you command? & Iron Cadence (military) & E & 4/4 & Yes & S4 & 3/3 \\
C7 & What do strangers actually share? & The Unnamed (no-stakes) & E & 4/4 & Yes & S4 & 3/3 \\
\midrule
\multicolumn{5}{l}{\textit{Aggregates}} & 7/7 (100\%) & 28/28 (100\%) & 7/7 & 7/7 (21/21) \\
\bottomrule
\end{tabular}
\end{table*}

\subsection{Completion Rate}

All 7 campaigns completed within 7 sessions. Completion rate: 100\%.

This is not a trivial result. Practitioner surveys estimate the completion rate
for human-run long-form TRPG campaigns at 30--50\%~\cite{perkins2017dungeon}. The
Marathon spine's 100\% rate reflects both the AI simulation context (no scheduling
pressures, no player attrition, no DM burnout) and the spine's structural design
(branching end-conditions ensure completion at any engagement level). These
contributions cannot be fully separated, but the structural design is necessary:
without explicit end-conditions, an AI-simulated campaign has no self-termination
mechanism. The spine is what makes the campaign able to end.

\subsection{External Completion Indicators}

The External Completion Indicators column in Table~\ref{tab:campaign-outcomes}
reports three rubric-external binary checks for each campaign (see
Section~\ref{subsec:operationalizing-completion} for definitions). All 7 campaigns
satisfy all 3 indicators: the PC-authored finale deliverable was produced and logged
verbatim in S7; all 4 hints were delivered before the finale; and the party achieved
Route D or E (the two highest-engagement end-conditions), or the campaign's designated
optimal route (Route B for C5, where Route B is the spine's highest achievable
end-condition for the professional-code archetype).

The alignment between rubric-based gate scores $\geq$75 and external completion
indicators (7/7 campaigns, 3/3 indicators each) suggests the gate score is a valid
proxy for campaign completion, though it is not identical to it. The gate score
captures session-by-session narrative quality across 8 dimensions; the external
indicators capture only the spine's three final-state outcomes. A campaign could
in principle satisfy external indicators while producing low gate scores throughout
(a well-designed finale arriving after thin sessions), or produce high gate scores
without satisfying external indicators (rich mid-campaign sessions without a finale).
The 7/7 alignment indicates neither scenario obtained in this corpus.

\subsection{Route Distribution and Archetype Comparison}

Table~\ref{tab:route-distribution} shows routes achieved across the 7 campaigns,
grouped by archetype.

\begin{table}[t]
\caption{Route Distribution by Campaign and Archetype}
\label{tab:route-distribution}
\begin{tabular}{@{}lllp{3.2cm}@{}}
\toprule
Campaign & Archetype & Route & Notes \\
\midrule
C1 & Grief-themed & D & Artifact-decision tree; 6 artifacts \\
C2 & Covenant & D & Shard recovery; 6 shards \\
C3 & Faith/authority & D & Dispensation evidence; keystone \\
C4 & Siege-resource & E & All 7 watch-tokens + supply die \\
C5 & Professional-code & B & False dossier; guild-marker system \\
C6 & Duty-spine & E & Unit cohesion die; field journal \\
C7 & No-designed-stakes & E & Reciprocal acts log \\
\midrule
\multicolumn{2}{l}{Route D or E} & 5/7 (71\%) & \\
\bottomrule
\end{tabular}
\end{table}

Five of 7 campaigns (71\%) achieved Route D or E. No campaign achieved only
Route A or C (minimum completion). Route B (C5, Thieves Guild) requires
clarification: for the professional-code archetype, the spine designated Route B
as the optimal achievable route, not an inferior result. The guild-marker system
(trust tokens earned through professional acts, not personal-stakes acts) was
designed with a lower route ceiling because the archetype's emotional register
is lower-affect by design. Route B in C5 is as complete an ending as Route D
in C1 --- they are different answers to different central questions at different
emotional registers.

The route distribution shows genuine differentiation: three campaigns at D,
three at E, one at B. This is strong evidence that the token/route system is
neither trivially achievable (which would produce all-D or all-E) nor impossible
(which would produce all-A). The spine produces calibrated outcomes that
vary with the campaign's archetype and the party's level of engagement with
the specific central question.

\subsubsection{The Strangers Result}

Campaign 7 (Party of Strangers, no designed personal stakes, Route E) is the
corpus's strongest archetype-independence datum. The strangers archetype was
designed as a structural stress test: no grief hooks, no professional codes,
no loyalty objects, no institutional identities. The PCs were people with nothing
in common, brought together by circumstance alone.

The spine produced Route E completion. The Reciprocal Acts log --- a lightweight
trust-accumulation mechanism designed specifically for the strangers context ---
functioned as a sufficient analogue to the grief-themed token systems of earlier
campaigns. The central question (\textit{what do people with nothing in common
actually share?}) was answered in the finale: Vess within sight of fire; Sera
who had seen them safe; Brek who had followed the wrong fork; Tam who sat beside him.

This result validates the spine's archetype independence at the structural level.
The five components work with any motivational content; the specific content of
each component must be calibrated to the archetype, but the structure generalizes.

\subsection{Hint Delivery Analysis}

Table~\ref{tab:hints} shows hint delivery by campaign and the number of campaigns
in which each hint required a recovery path.

\begin{table}[t]
\caption{Hint Delivery Rate by Campaign (4 Hints per Campaign)}
\label{tab:hints}
\begin{tabular}{@{}llr@{}}
\toprule
Campaign & Delivery status & Recovery paths used \\
\midrule
C1 Moon-Silver & 4/4 (100\%) & 0 \\
C2 Conclave Compact & 4/4 (100\%) & 1 (H3 extended delivery) \\
C3 Halted Spire & 4/4 (100\%) & 1 (H2 recovery active S2) \\
C4 Thorngate Watch & 4/4 (100\%) & 0 \\
C5 Thieves Guild & 4/4 (100\%) & 0 \\
C6 Military Unit & 4/4 (100\%) & 0 \\
C7 Strangers & 4/4 (100\%) & 1 (H1 passive-extended) \\
\midrule
Total & 28/28 (100\%) & 3 campaigns \\
\bottomrule
\end{tabular}
\end{table}

All 28 hints across 7 campaigns were delivered (100\%). Three campaigns used
documented recovery paths; in each case, the recovery path activated cleanly
because it was specified in the spine document before the campaign began. The
recovery-path mechanism was critical for these 3 campaigns: without a pre-specified
secondary delivery mechanism, the missed primary delivery would have created a
permanent gap in the hint structure.

Hint delivery timing across the corpus showed a consistent pattern:
Hint 1 was delivered in Session 1 or 2 in all 7 campaigns.
Hint 2 was delivered in Sessions 2--4 in all 7.
Hint 3 was delivered in Sessions 4--6 in all 7, with the C2 recovery path
extending H3 delivery into Session 5 from the designed Session 4 delivery.
Hint 4 was delivered in the finale session (Session 7 or late Session 6) in
all 7 campaigns, in most cases via auto-delivery at the designed passive threshold.

The simultaneous convergence event --- all four hints available to the party's
recognition at the finale moment --- occurred in all 7 campaigns. In no campaign
was the convergence scripted by the DM; in each case, it emerged from the party's
arrival at the finale scene with all four hints previously delivered.

\subsection{Sessions-to-First-Binding-PASS}

Mean sessions-to-first-binding-PASS across all 7 campaigns is 4.5 (range: 4--6).
Six of 7 campaigns achieved first PASS at Session 4; one (C4) achieved first PASS
at Session 6, reflecting the siege-resource archetype's extended trust-building period.

No campaign needed all 7 sessions to achieve first binding-mode PASS. The
advisory period (Sessions 1--3) is structurally appropriate: it allows the campaign
to establish its material, its character relationships, and its hint-delivery
trajectory before imposing binding-quality thresholds. The data supports the
advisory-period design as a campaign architecture decision, not merely a
rubric-scoring choice.

The 4.5 mean is archetype-independent. The grief-themed campaigns (C1--C3) achieve
first PASS at S4--S5 without special structural features; the professional-code
campaign (C5) achieves first PASS at S4; the no-stakes campaign (C7) achieves
first PASS at S4 despite the absence of personal-stakes material. The spine's
structure brings sessions to binding quality reliably within the first two-thirds
of the arc regardless of archetype.

%% sections/05-discussion.tex

\section{Discussion}
\label{sec:discussion}

\subsection{Why the Spine Works: Structural Analysis}

The campaign spine's 100\% completion rate and 71\% Route D/E rate cannot be
attributed to any single component. The five components function as a system:
the central question gives the hints their convergence target; the hints give
the routes their achievability conditions; the spotlight schedule gives the
deliverable its raw material; and the deliverable gives the finale its
non-scripted emotional content. Remove one component and the system degrades;
the corpus shows no campaigns designed with partial spine implementations,
but the component-interaction analysis in Section~\ref{sec:components}
predicts specific failure modes for each absent component.

The most critical component for reliable \textit{completion} is the branching
end-conditions. Campaigns without explicit completion criteria cannot
self-terminate, and in AI-simulated play this produces indefinite continuation
rather than a clean ending. The route structure guarantees that the campaign
can end gracefully at any engagement level; Routes A--C provide the structural
floor.

The most critical component for \textit{quality} of completion --- as measured
by Route D/E achievement and deliverable richness --- is the four-hint schedule.
Without Hint 3 and Hint 4, the party has no basis for recognizing that the
optimal ending is available. The hints do not force the party to reach Route D/E;
they create the structural possibility for it. A party that reaches the finale
with all four hints already delivered can achieve the optimal route; a party
that reaches the finale missing Hint 3 cannot, regardless of their intent
or engagement level.

\subsection{The Advisory Period as Campaign Architecture}

The advisory-to-binding calibration --- Sessions 1--3 as advisory, Session 4
onward as binding --- is usually framed as a rubric-scoring decision (scores
in advisory sessions are not held to binding-mode thresholds). The evidence
supports reframing it as a campaign design decision. Sessions 1--3 are the
hint-planting window for Hints 1 and 2; they are the sessions in which PCs
encounter their spotlight material for the first time; they are the sessions
in which trust-accumulation systems are first loaded. Sessions 4--7 are the
sessions in which that material is used. A campaign that attempted binding-mode
quality in Session 1 would be structurally premature: the material for
binding-quality sessions depends on the planting work of Sessions 1--3.

The 4.5 mean sessions-to-first-binding-PASS validates this: in most campaigns,
the advisory period provides exactly the preparation the binding sessions require.
The exception (C4, first PASS at Session 6) is consistent with C4's extended
trust-building structure (7 watch-tokens, supply die attrition, HP carry-over):
C4 required more preparation sessions because its trust-mechanism was more
demanding.

\subsection{Route B as Designed Optimal: The C5 Case}

The professional-code archetype (C5, Thieves Guild) presents a result that
requires careful interpretation. The campaign achieved Route B, not D or E,
and Route B is listed in the route taxonomy as a ``personal-level resolution''
outcome rather than a ``full inversion'' outcome. Does this represent a spine
failure?

The answer is no, and the distinction matters. The C5 spine was designed with
Route B as the campaign's optimal achievable route. The guild-marker system
(trust tokens earned through professional acts, not personal-stakes acts)
was calibrated to a lower emotional register than the grief-themed or
covenant-themed campaigns. The central question (\textit{what does professional
loyalty cost?}) does not have a ``full inversion'' answer; it has a
``professional resolution'' answer, which is Route B. A Route D outcome for
C5 would have required personal-stakes material the archetype was explicitly
designed to exclude.

This indicates that the route-naming convention (A--E as difficulty-ordered)
should not be read as quality-ordered. Route E in a siege-resource campaign is
not inherently better than Route B in a professional-code campaign; they are
different kinds of answers to different kinds of questions. The spine's design
requirement is that the campaign's designated optimal route is achieved, not that
the letter designation is as high as possible.

\subsection{The Strangers Result and Archetype Independence}

The no-designed-stakes archetype (C7) achieving Route E is the corpus's most
consequential archetype-independence result. Earlier hypotheses about the spine's
scope held that the grief-hook structure and personal-stakes material in PCs'
character sheets were prerequisites for Route D/E completion. Campaign 7
falsifies this. The Reciprocal Acts log --- a system that records specific acts
of consideration between strangers without requiring grief-hooks or professional
codes --- produced Route E completion with an archetype that had neither.

The implication for the spine's generalizability is significant: the five
components do not require emotionally dense PC material to function. They require
\textit{some} form of trust-accumulation system calibrated to the archetype,
but the system's form is flexible. Grief-tokens and professional-markers and
reciprocal-acts are all instances of the same structural mechanism: specific
party actions earn specific evidence of engagement with the central question,
and that evidence gates the route conditions.

This suggests the spine is exportable to TRPG systems and settings with lower
personal-stakes norms than Dragonlance. Whether it functions equally well with
purely tactical parties (no social material at all) remains an open question;
it is the obvious next stress test.

\subsection{Limits of the Current Corpus}

Three limitations bound the current claims.

\textbf{AI simulation confound.} The 100\% completion rate is confounded by the
AI simulation context. No player attrition, no scheduling failures, no DM burnout.
Human-run campaigns using the spine would face all three. The spine's structural
design is a necessary but not sufficient condition for completion in human play;
this paper establishes sufficiency in the AI-simulated context only.

\textbf{Seven-session constraint.} All campaigns were designed for exactly 7 sessions.
This prevents generalization to shorter (5-session) or longer (10-session) campaigns.
The 7-session length was chosen as sufficient for full 4-hint delivery, 4-PC spotlight
distribution, and Route D/E conditions. Whether the spine's components scale to different
arc lengths is unexamined. In particular, a 5-session spine would need to compress the
hint delivery schedule (Hints 1 and 2 in Session 1; Hints 3 and 4 in Session 4), which
may reduce hint differentiation and convergence impact.

\textbf{Single designer.} All 7 campaign spines were designed by the same
researcher-designer. Cross-designer reliability --- whether other designers
implement the five components with equivalent outcomes --- is not established.
The component specifications are designed to be implementable from written
guidance, but this has not been tested with naive implementers.

\subsection{Future Work}

Three research directions follow directly from the corpus:

\textbf{Human-run validation.} The most important extension is running the spine
in human-DM tabletop play. The five components translate directly to human play;
the recovery-path mechanism is especially important in that context (human DMs
miss delivery windows more frequently than AI DMs following explicit condition
specifications). A human-run corpus of even 3--4 campaigns would substantially
strengthen the generalizability claim.

\textbf{5- and 10-session variants.} Testing the spine at arc lengths other than
7 would establish its flexibility bounds. The prediction is that 5-session arcs
will succeed with a compressed hint schedule but will produce lower Route D/E
rates; 10-session arcs will produce equivalent or higher Route D/E rates but risk
drift in sessions 6--9 (the window between Hint 4 delivery and the finale).

\textbf{Partial spine implementations.} The corpus contains no campaigns designed
with partial spine implementations (e.g., hints and routes but no deliverable
specification; deliverable and spotlights but no explicit end-conditions). A
controlled study varying which components are included would directly test the
component contribution analysis. The prediction --- that branching end-conditions
and hint schedule are necessary while spotlights and deliverable specification
are quality-enhancing rather than completion-necessary --- is testable.

The circularity concern raised by our operational definition of completion is
partially addressed by the three external indicators described in
Section~\ref{subsec:operationalizing-completion}: deliverable produced, hints
delivered, and Route D/E achieved. These indicators are observable without
rubric scoring and provide convergent-validity evidence for the Route D/E
classification. However, a fully independent measure of ``emotional completeness''
--- such as player-written post-campaign reflections documenting their
experience of the finale, or narrative coherence scoring by readers unfamiliar
with the rubric who assess whether the campaign's central question was answered ---
remains future work. Such independent measures would allow the rubric-based gate
score and the external-indicator composite to be validated against a criterion
that shares no measurement ancestry with either.

%% sections/06-conclusion.tex

\section{Conclusion}
\label{sec:conclusion}

Multi-session TRPG campaign arcs fail to complete at high rates, and when they
complete, they often do so through DM improvisation rather than designed architecture.
The emotional payoffs that make long-form TRPG play valuable --- the arcs that land,
the finales that feel earned, the moments players cite years later --- are structurally
unreliable in the absence of explicit campaign design.

The campaign spine addresses this unreliability through five components that together
specify what the campaign is about (central question), how the party will come to
understand it (hint schedule), what constitutes completion at different engagement
levels (branching end-conditions), when each character's arc is foregrounded
(spotlight schedule), and what the finale's emotional content will come from
(PC-authored deliverable). Across 7 complete campaigns in the Marathon AI-simulated
TRPG workshop, these five components produced:

\begin{itemize}
  \item \textbf{7/7 campaign completions} (100\% completion rate; 49 sessions total)
  \item \textbf{5/7 Route D or E achievements} (71\%; the end-conditions requiring
    highest sustained party engagement)
  \item \textbf{28/28 hints delivered} (100\%; three campaigns used recovery paths)
  \item \textbf{7/7 PC-authored deliverables produced} verbatim at the finale
  \item \textbf{Mean sessions-to-first-binding-PASS: 4.5} (range: 4--6 across
    all campaign archetypes)
\end{itemize}

The spine is archetype-independent. It produced equivalent completion and
quality outcomes across grief-themed, covenant, faith/authority, siege-resource,
professional-code, duty-spine, and no-designed-stakes campaign types. The
strangers result (Campaign 7, Route E, no designed personal stakes) is the
strongest evidence for this claim: the spine's five components produced Route E
completion for a party with no grief-hooks, no professional codes, and no
institutional identities, through a lightweight trust-accumulation mechanism
calibrated to the archetype's emotional register.

The paper's primary contribution is not the individual components --- hint
scheduling, route branching, and PC spotlights appear in various forms in
practitioner TRPG design resources. The contribution is the demonstrated
\textit{system}: five components that interact specifically and predictably,
producing reliable emotional completeness across diverse campaign archetypes
when all five are implemented together. The interaction structure
(Table~\ref{tab:component-interaction}) is the design knowledge this paper
contributes: hints enable routes; spotlights load deliverables; end-conditions
determine hint delivery priority; central questions drive hint coherence.

We make the full campaign spine template available as supplementary material,
including the complete specifications for all 7 campaigns in the corpus. Our
hope is that campaign designers --- whether running AI-simulated or human-run
TRPG --- will find the five-component structure useful as a design constraint
that does not determine narrative content but provides the architectural
backbone without which narrative content cannot reliably converge into an
emotionally complete arc.


\bibliographystyle{ACM-Reference-Format}
\bibliography{references}

\end{document}
